\chapter*{Conferencia XIV}
\addcontentsline{toc}{chapter}{LECTURA I - El timbre y el carácter del violín}
SEÑORES: Ahora voy a presentarles la máxima uniformidad del sonido del violín. Confieso sentir bastante temor al abordar este problema. Quizás el temor me lleva a posponer su presentación hasta casi el final del curso. Quizás el hecho de que mi solución para este problema solo haya recibido una demostración contribuye a aumentar mi temor. Sé muy bien que una golondrina no hace verano. También sé muy bien que los autores más engañosos sobre el violín son aquellos que describen las peculiaridades del sonido de un solo violín. Presentar las peculiaridades del sonido de un solo violín como hechos que existen en todos los violines debería hacer que el autor sea susceptible de ser procesado penalmente por conspiración para defraudar.
Solo los factores que dependen de la acción del aire pueden considerarse constantes en la influencia del tono del violín; mientras que todos los factores que dependen de la acción de la madera no son fiables ni constantes, sino que, por el contrario, son caprichosos en su acción de principio a fin. Es esta acción caprichosa la que impide la aplicación de reglas estrictas a la construcción de violines. Es evidente que el mercado estaría saturado de violines "de primera calidad" si se aplicaran reglas estrictas a su construcción. Es la acción caprichosa de la madera la que hace que el violín...
pertenece a una clase aparte. Es perfectamente posible regular el tono de todos los demás instrumentos musicales con tal uniformidad que el tono de uno sea exactamente igual al de otro; mientras que ningún grado de habilidad humana puede hacer que dos violines suenen exactamente igual.
205

similares salvo en raras ocasiones.
Por lo tanto, quienquiera que describa las características de un solo violín como infalibles en cuanto a su sonido, es culpable de fraude, ya sea intencional o no. Reitero que solo los modificadores del sonido del violín que dependen de la acción del aire pueden considerarse factores infalibles.
Teniendo esto presente, les pido que también lo tengan en cuenta al estudiar o aplicar mis detalles para lograr la máxima uniformidad en el sonido del violín. Que estos detalles puedan fallar a veces y tener éxito otras es tan seguro como la acción caprichosa de la madera. Sin ningún deseo de engrandecer mi propio logro,
Afirmo que estos detalles para la uniformidad de la potencia tonal del violín fueron trazados en papel como resultado de un razonamiento a priori. Este hecho me humilla. El hecho de haber dedicado toda una vida al estudio del timbre del violín y solo haber logrado obtener un único factor beneficioso mediante razonamiento a priori es humillante. Sin embargo, de ninguna autoridad puedo encontrar otro factor beneficioso para el timbre del violín que haya sido trazado en papel con anterioridad. A lo largo de los 400 años de historia del violín, ningún hecho es tan prominente como el hecho de que cada factor beneficioso para el timbre del violín se debe directamente a la experimentación. Hasta ahora, el trazado en papel de cada uno de estos factores ha sido resultado de un razonamiento a posteriori. Prácticamente, para el violín no hay diferencia a favor de ninguno de los dos métodos para trazar factores beneficiosos. El único beneficio reside en el estudiante de violín; y dicho beneficio consiste enteramente en el estímulo. Si
206

Si el pensamiento puede asegurar un factor que beneficie el sonido del violín, entonces puede asegurar dos factores de ese tipo. ¡Ahí está la clave! Decir que "el violín alcanzó la perfección hace 200 años" es profundamente desalentador para quienes lo aceptan como cierto. Afortunadamente, ahora se entiende generalmente que este dicho, en los últimos años, se originó con el promotor del comercio de violines antiguos.
La aplicación de los detalles para lograr la máxima uniformidad de la potencia tonal se realizó en un solo violín; desde entonces, esta silla de ruedas de tres ruedas, desde la cual me dirijo a ustedes, ha ocupado mi atención. Tan marcada fue la uniformidad de la potencia tonal en esta única demostración que me siento animado a pedirles que retomen mi solución donde la dejo; y, si mi solución resulta correcta, entonces estaré ampliamente agradecido.
recompensado.
Sin diagramas, me enfrento a una dificultad inusual para seleccionar las palabras que describan con precisión los detalles y logren la máxima uniformidad en la intensidad del tono; por lo tanto, es posible que necesiten reflexionar detenidamente para comprender su descripción. En tal caso, no conozco otra manera que repasar el tema repetidamente. Esta repetición requiere esfuerzo, pero les brindo mi apoyo.
Como medio para ayudar a resaltar los detalles de la uniformidad de la potencia tonal con claridad, primero es necesario presentar las líneas de razonamiento que llevaron a ese hecho. Solo con pasos lentos y pensamiento concentrado pude seguir esas líneas hasta una conclusión; e incluso procediendo así,
207

Se encontró la necesidad de salvar un abismo enorme con una simple "suposición". Como habrás observado en situaciones de emergencia, "suponer" a menudo se convierte en una
necesidad.
Cada instrumento musical posee un tono fundamental, o más bajo, propio de sí mismo; y, en todos los instrumentos musicales, excepto en los de cuerda, se prevé la producción de tonos más agudos que el tono fundamental. En todos los instrumentos de viento, dicha previsión se basa enteramente en la acción del aire; mientras que, en los de cuerda, dicha previsión se basa en parte en la acción del aire y en parte en la del viento.
la acción de la madera y las cuerdas. Esta diferencia
coloca a la familia de cuerdas en una clase aparte; además, esta diferencia complica enormemente los problemas relacionados con la producción de tonos de cuerda; además, esta diferencia desafía la habilidad humana para producir una uniformidad precisa en los valores de los tonos de cuerda. Aunque el piano es un dispositivo musical de cuerda,
Sin embargo, debido a que cada uno de sus tonos es un tono fundamental, mientras que el violín solo tiene cuatro tonos fundamentales, por lo tanto existe una relación remota entre estos dos instrumentos. Porque el piano
El tono depende del golpe directo de un martillo; por lo tanto, el piano pertenece a la clase de instrumentos de percusión, a pesar del uso de una caja de resonancia para aumentar la potencia del sonido. El tono del violín depende del enrollado y desenrollado de las cuerdas y del uso de una caja de resonancia para aumentar la potencia. Aquí termina la similitud entre estos dos instrumentos. La caja de resonancia del piano debe ser alargada,
208

acortado, ensanchado y regulado en grosor para acomodar todos y cada uno de sus tonos fundamentales, mientras que esa parte particular de la caja de resonancia del violín, que aumenta el tono fundamental de cada cuerda, también debe aumentar todos los demás posibles
tonos en cada cuerda. Esta disimilitud permite la uniformidad de la potencia del sonido en el piano y provoca la falta de uniformidad de la potencia del sonido en el
violín.
La reflexión y la experimentación, centradas en la escala del piano (longitud comparativa de las cuerdas y la caja de resonancia), han dado como resultado una uniformidad de sonido bastante satisfactoria. Sin embargo, la caja de resonancia del violín aún presenta deficiencias en esta calidad de sonido. Naturalmente, surge la pregunta: "¿Pueden la reflexión y la experimentación mejorar la caja de resonancia del violín para lograr una mayor uniformidad en el sonido?".
Para resolver esta cuestión, primero es necesario señalar claramente los errores en las dimensiones de la caja de resonancia del violín. Naturalmente, quienes creen que el violín alcanzó la perfección hace 200 años se negarán a admitir que existen imperfecciones en la caja de resonancia de hace 200 años; o que existen imperfecciones hoy en día en cajas de resonancia precisamente similares a la de hace 200 años.
Posiblemente, ninguna evidencia alguna pueda cambiar tal cosa.
creencia.
Como método exitoso para mantener el error, no hay ninguno tan exitoso como el rechazo de* evidencia. Al recurrir a la caja de resonancia del piano para obtener evidencia de error en la caja de resonancia del violín,
209

Seguramente habrá quienes se nieguen a recibir dichas pruebas por temor a ser convencidos en contra de su voluntad.
A continuación se presentan dichas pruebas:
(1) La caja de resonancia del piano es más pesada debajo de las cuerdas más gruesas.
(2) La caja de resonancia del piano es más ligera debajo de las cuerdas más pequeñas.
(3) La caja de resonancia del piano es más larga debajo de las cuerdas más gruesas.
(4) La caja de resonancia del piano es más corta debajo de las cuerdas más pequeñas.
En la caja de resonancia del Stradivarius no hay diferencia en la longitud de la superficie de contacto con las cuerdas. Tampoco hay diferencia en el grosor de dicha superficie. A partir de estos datos, resulta evidente que la caja de resonancia del violín de hace 200 años no había alcanzado la perfección. [En cuanto a la uniformidad del sonido, algunas cajas de resonancia de Joseph Guarnerius se acercaron mucho más a la perfección que cualquier caja de resonancia de Stradivarius.]
Debido a que la caja de resonancia del piano actual produce una mayor uniformidad en la potencia del tono, es evidente que el error reside en la caja de resonancia del violín; y dicho error se hace evidente al intentar producir dos octavas de tono, con potencia uniforme, en cualquiera de las cuatro cuerdas del violín. Incluso con la ayuda de una mayor presión del arco, tal intento resulta infructuoso. La razón de este fracaso es evidente; la misma longitud y grosor de la caja de resonancia, dedicada a aumentar el tono fundamental,
210

También debe emplearse para aumentar el tono de la cuerda acortada y, por consiguiente, debilitada. De este hecho, solo hay una conclusión: que tanto la longitud de la actividad de la caja de resonancia como su rigidez son excesivas para el golpe debilitado de la cuerda acortada. Partiendo del tono fundamental de cualquier cuerda de violín y contando los semiintervalos, hay veinticinco tonos dentro de dos octavas que deben ser aumentados por un conjunto idéntico de fibras de la caja de resonancia. En el piano, cada uno de estos tonos sucesivos se aumenta mediante una caja de resonancia de menor longitud y menor grosor. En el violín, tal longitud y rigidez reducidas de la caja de resonancia, para cada tono sucesivo, son manifiestamente imposibles. De hecho, igualar la uniformidad de la potencia sonora del violín con la del piano sería una esperanza vana si no fuera por la ayuda de una mayor presión del arco.
Resulta evidente que lograr una potencia tonal uniforme en dos octavas en cada cuerda del violín es un problema plagado de dificultades. También resulta evidente que una solución a este problema es valiosa para el sonido del violín; y, en conciencia, pretender mejorar los excelentes métodos de construcción de la caja de resonancia de Stradivarius y Joseph Guarnerius es, al parecer, motivo suficiente para dudar. Durante 200 años, esos métodos han sido fielmente copiados por los constructores de violines más ambiciosos del mundo civilizado; y, en todo el mundo, no conozco a nadie, ni he oído hablar de nadie, que afirme haber mejorado los mejores métodos de esos dos famosos.
211

experimentalistas. Sin embargo, a pesar de que no se ha encontrado a nadie que haga tal afirmación, conozco violines modernos con un sonido tan bello como cualquier Stradivarius que haya escuchado, y, si bien tienen un tono igualmente dulce, poseen una mayor uniformidad en la potencia del sonido.
¿Por qué no reclamar lo que se debe?
¡No veo ningún inconveniente en reclamar lo que se debe cuando el reclamante tiene los bienes para demostrarlo!
Es cierto que lo mejor es no reclamar más de lo que se merece, y también lo mejor es ser modesto al hacer cualquier afirmación sobre uno mismo; pero de nada sirve ocultar el propio talento. El mundo siempre está dispuesto a reconocer el mérito cuando se convence de que es merecido; pero, con razón, el mundo exige pruebas antes de juzgar.
Considerando todos los aspectos, presente la prueba.
Al presentar tales pruebas, no hay a quién temer excepto al promotor del comercio de violines antiguos; y, ahora se entiende de forma bastante generalizada que está destinado a ese puerto donde se envían los mentirosos deliberados.
Como recordamos vívidamente los estudiantes de violín de mayor edad, Ole Bull podía extraer cuatro octavas de un tono musical bastante uniforme de cada cuerda de su "Joseph". Como saben todos los estudiantes de violín, es una falacia afirmar que tal mérito reside enteramente en el intérprete. No importa quién sea el intérprete, primero debe tener el instrumento en sus manos. *Es mi propia observación que el "Joseph", entre los violines antiguos, sigue siendo inigualable como violín solista en los auditorios más grandes; y esta opinión se basa en el hecho de que el "Joseph" posee mayor potencia y uniformidad de tono. Es natural que
212

El estudiante pregunta: "¿Qué causa mayor potencia y mayor uniformidad de la potencia tonal en el 'José'?" Durante muchos años, he intentado responder a esta pregunta, pero, por la misma razón que se deja a la conjetura, los datos precisos para el "José" se conservan como un valioso activo personal. La naturaleza humana sigue siendo muy humana. En cuanto a los datos para el "José",
Honeyman afirma lo siguiente: "El 'vientre' es más grueso
en los bordes y más delgada en las zonas centrales." De una descripción tan imprecisa, poco se puede aprender. Incluso un autor tan bien preparado por su educación, por su formación exhaustiva en la construcción de violines y por la oportunidad de observar como Ed. Herron Allen, remite a sus lectores a los datos de Stradivarius sobre el grosor de las tablas del "Joseph", aunque afirma haber tenido el violín Joseph de M. Sainton para obtener los datos.
¡Extraño!
He aquí otro caso aún más extraño. Envié a París el libro y los gráficos del Sr. Simoutre porque entre ellos había uno de José. Recibí el libro con el gráfico de José recortado.
"¡Diable!"
De un amigo obtuve una copia de un diagrama del violín "José". El original de este diagrama fue realizado en el taller de Vuillaume. A partir del diagrama que recibí y de la descripción imprecisa que Honeyman hizo del violín José, dediqué mucho tiempo a realizar repetidos experimentos. Hoy no lamento la pérdida del diagrama de Simoutre, y, debido a la opinión de que es imposible dedicar más tiempo a experimentar con las posibles variedades de violín,
213

El grosor de la placa es mayor que el que he indicado. Los gráficos para el Stradivarius, de varias variedades diferentes en grosor de placa, se obtienen con bastante facilidad. No solo he realizado repetidas pruebas con todas las variedades de Stradivarius en grosor de placa, sino que también he realizado muchas pruebas con los datos disponibles para el Joseph Guarnerius. Entre estos dos grandes genios, sin dudarlo, atribuyo el mérito de una mayor potencia tonal y una mayor uniformidad de la potencia tonal a Joseph Guarnerius. Por la mayor duración del tono y por la mayor dulzura del tono, atribuyo el mérito a Stradivarius.
El método de medición del espesor de la placa, finalmente adoptado por estos dos experimentadores exitosos, se presenta con el propósito de resaltar con claridad mis líneas de razonamiento, que conducen a un método parcialmente nuevo.
Así pues: En los experimentos realizados con el método Joseph, se demostró claramente que se obtenía una mayor potencia tonal al colocar el punto más delgado de la caja de resonancia a medio camino entre la posición del puente y los extremos de la placa.
Obsérvense los dos hechos siguientes en los métodos de Strad y Joseph:
(1) El punto más delgado debajo de todas las cuerdas está a igual distancia de la posición del puente, independientemente de la diferencia en el peso y el diámetro de las cuerdas.
(2) Igual espesor de la placa debajo de todas las cuerdas, independientemente de la diferencia de peso y diámetro de las cuerdas.
En este momento solo llamo brevemente la atención sobre las diferencias entre este método para el grosor de la caja de resonancia y el método (mi método) para
214

aún mayor uniformidad en la potencia del sonido. El método favorito de Stradivari es:
(1) El punto más delgado de la caja de resonancia se encuentra a la mayor distancia práctica del puente para todas las cuerdas, sin tener en cuenta las diferencias de peso y diámetro de las cuerdas.
(2) El grosor de la caja de resonancia es igual debajo de todas las cuerdas sin tener en cuenta la diferencia de peso y diámetro de las cuerdas.
Comparando el nuevo método, de la siguiente manera:
(1) El punto más delgado debajo de la cuerda G está a la mayor distancia práctica de la posición del puente.
(2) El punto más delgado debajo de la cuerda D está menos distante que para la cuerda G.
(3) El punto más delgado debajo de la cuerda A está menos distante que para la cuerda D.
(4) El punto más delgado debajo de la cuerda E está menos distante que para la A.
(5) En la posición del puente, el mayor grosor de la caja de resonancia se encuentra debajo de la G.
(6) En la posición del puente, menor grosor de la caja de resonancia para D que para G.
(7) En la posición del puente, menor grosor de la caja de resonancia para A que para D.
(8) En la posición del puente, el grosor de la caja de resonancia es menor para E que para A.
Este nuevo método se basa tanto en hechos como en
teoría.
Lo cierto es que las cuerdas varían en peso y diámetro; por lo tanto, la fuerza que ejercen al accionarlas también varía.
La teoría es: El grosor de la caja de resonancia debajo de cada una debe variar según el diámetro y el peso de
215

cuerdas; y la duración de la actividad de la caja de resonancia debajo de cada cuerda debe variar según el tono de los sonidos que se exigen de cada cuerda.
Para este hecho, no se necesita prueba alguna; es evidente por sí mismo; pero la teoría sí necesita prueba. Como prueba, primero presento como evidencia, y una demostración de este hecho, que cada cuerda depende en gran medida (aunque no totalmente) de la actividad de la caja de resonancia justo debajo para la amplificación del sonido. Para tal demostración, ofrezco este violín como sacrificio. Es cierto que solo ofrezco uno para el beneficio de muchos, pero, de alguna manera, me siento como un carnicero, es decir, le concedo cualquier sentimiento al carnicero. Con esta delgada hoja, procedo a partir la caja de resonancia de este violín en varios lugares, y continuaré tal trabajo hasta que el sonido de todas las cuerdas esté completamente arruinado.
(1) Comenzando en el extremo inferior de la salida derecha, dividí la caja de resonancia hasta la purga. La aplicación del arco no muestra ningún daño al tono de ninguna cuerda como sigue a esto
dividir.
(2) Comenzando en un punto cercano y debajo del poste, partí la caja de resonancia hasta el bloque del extremo inferior. Al aplicar el arco nuevamente, no se observa ningún daño al tono en ninguna cuerda.
(3) Comenzando cerca del puente y debajo de él, partí la caja de resonancia a lo largo de la junta central hasta el bloque del extremo inferior. Nuevamente, no hubo daño al tono de ninguna cuerda.
(4) Comenzando en el extremo inferior de la salida izquierda, divido la caja de resonancia hasta el
216

I
purng. De nuevo, sin daños en el tono.
(5) Comenzando en el punto a la altura del puente, y a la izquierda de la barra, partí la caja de resonancia hasta el bloque inferior. Aunque no se podría llamar ruina total, ahora tanto el tono G como el D muestran daños graves.
(6) Comenzando debajo del puente y a la derecha de la barra, partí la caja de resonancia hasta el bloque inferior del extremo. El daño a los tonos G y D aumenta, pero aún no aparece daño al tono en A y E.
(7) Comenzando en un punto a la altura del puente, y una pulgada a la izquierda, partí la caja de resonancia hacia arriba hasta el púlpito. A partir de esta división, solo hay un ligero daño adicional al tono G, pero no hay daño adicional al tono D.
(8) Comenzando en el puente, cerca y a la izquierda de la barra, partí la caja de resonancia hacia arriba hasta el bloque final. El tono G ahora está completamente arruinado, mientras que el tono D, aunque dañado, aún no está totalmente arruinado.
(9) Comenzando en el puente, y debajo del D, partí la caja de resonancia hacia arriba hasta el bloque final. El tono D ahora está completamente arruinado, mientras que el tono A, aunque dañado, no está totalmente arruinado. (10) Como antes, partiendo debajo de A completamente
arruina el tono de esa cuerda, mientras que el tono de Mi solo se arruina parcialmente. (11) Como antes, la división debajo del Mi arruina completamente el tono de esa cuerda. Así se demuestran claramente las áreas de actividad de la caja de resonancia de las que depende cada cuerda para el aumento de la potencia del tono; y sus definitivas
217

La ubicación demostró ser de inmenso valor para la producción de una potencia tonal uniforme. Es evidente que
Dicha ubicación de las áreas permite una reducción precisa del grosor de la caja de resonancia debajo de cada cuerda, proporcional al diámetro y al peso de cada cuerda.
Hasta ahora, la evidencia es concluyente; pero el siguiente paso es un camino incierto. La pregunta es: "¿Cuál será la proporción de longitudes para la actividad de la caja de resonancia que mejor potencie los tonos de cada cuerda?"
La proporción para acortar la caja de resonancia del piano no se puede aplicar al violín; ni en toda la gama de instrumentos musicales puedo encontrar una proporción que se pueda aplicar al violín.
Al hojear un libro de texto dedicado a la filosofía del sonido musical, encontré una proporción para las quintas de la escala mayor. Esta proporción atrajo inmediatamente mi atención y, debido a que no solo señala la diferencia en el número de vibraciones por segundo para dos o tres quintas consecutivas,
pero, para todas las quintas posibles en las escalas mayores.	Él
es una razón constante, ese es el tipo de razón que yo	
ahora lo deseo.	El libro afirma que multiplicar
multiplicando el número de vibraciones de cualquier tono por 3-2, se obtiene el número de vibraciones en su quinta superior. Deseando certeza en este asunto, procedo a comprobar la constancia de esta proporción aparentemente inocente, 3-2. Helmholtz afirma que conoces a Helmholtz, debe ser alemán, creo que de todos modos, porque nos dice más sobre la filosofía del sonido musical que todos los demás filósofos de todos los demás países juntos. Helmholtz afirma que el sol abierto,
218

Un violín afinado en tono de concierto vibra 200 veces por segundo; por lo tanto, si se multiplica 200 por esta sencilla proporción, 3-2, se obtiene la quinta por encima del sol al aire; 200, multiplicado por 3-2, es igual a 300; 300, multiplicado por 3-2, es igual a 450; 450, multiplicado por 3-2, es igual a 675. Estas cifras representan los tonos al aire del violín: sol, re, la, mi. ¡Qué fácil! Después de todo, este asunto de la filosofía no es tan complicado. Cualquiera puede hacerlo. Sí, si 3-2 funciona para la música, sin duda debería funcionar bien para el violín; sin embargo, recuerdo casos en los que no parecía haber ninguna proporción visible de ningún tamaño entre la música y el violín. Pero, como busco una proporción constante para las longitudes de actividad de la caja de resonancia bajo las cuerdas del violín, y como no parece haber otra proporción que no sea 3-2 en el violín, por lo tanto, estoy obligado a...
prueba 3-2.
Por lo tanto: Considerando que la mayor longitud de actividad de la caja de resonancia debajo de G es de 12 pulgadas, por lo tanto, todo lo que se necesita para encontrar dicha longitud para D es usar esta proporción aparentemente inocente, 3*2,.-- sobre , 12;
y, 12 multiplicado por 3-2 -es igual a 18 ¡qué!	18
pulgadas de actividad de caja de resonancia para D! Ojalá yo"	
Pude ver a Helmholtz por un minuto.	-
En la vida existe algo llamado confianza mal depositada. A veces se deposita en los demás; otras veces, en nosotros mismos. Esta última es la peor clase, no podemos insultar a Jones.
Sin embargo, a pesar de haber sido llevado por el 3-2, dediqué muchos meses a buscar una proporción que se aplicara a la duración de la actividad de la caja de resonancia bajo las cuerdas. Pero, tanto en las horas de vigilia como en las de sueño, dondequiera y cuando...
219

Siempre que busqué tal proporción, allí estaba ese 3-2 como el proverbial fantasma. No se levantaba, ni se movía, ni desaparecía; permanecía tan imperturbable como un indio de madera. Llegué a odiar el 3-2; pero, cuanto más lo odiaba, más lo tenía. Desesperado, busqué precedentes que dieran alivio a los fantasmas. Recordé tratamientos exitosos de otros que sufrían de "fantasmas", pero el mío no era de ese tipo. Una noche de insomnio, encontré el precedente que tanto había buscado. Era el precedente proporcionado por el capitán de una barcaza del Misisipi. Este capitán, al oír hablar de un invento novedoso capaz de predecir la llegada de tormentas, invirtió una buena suma en un barómetro y lo instaló cerca del timón. Poco después, una gran tormenta amenazó con hundir la barcaza. El barómetro, al no estar involucrado en tormentas, permaneció tranquilamente indiferente. Para el capitán, la diferencia entre una tormenta de viento y una tormenta eléctrica con truenos y relámpagos era insignificante. Sintiéndose víctima de una confianza mal depositada, el capitán se aferró a ese barómetro impasible y lo puso patas arriba. Este era mi precedente. Dado que ese capitán luego logró un empate exitoso, yo también me levanté, aproveché ese impasible 3-2 y lo puse patas arriba.
Luego dormí.
Confieso que esta nueva proporción, 2-3, es un mero capricho; pero el violín mismo es producto de caprichos. Durante sus primeros 200 años de desarrollo, el violín debe todo a caprichos; luego, debido a dos
220

Los afortunados caprichosos se convirtieron en genios, la suerte parece haber abandonado a la clase; es decir, si tan solo hubiéramos confiado en el viejo promotor de violines. Creo que, según el precedente de 200 años, ya es hora de tener más suerte.
